%% Lecture 14 -- Magnetostatics, Biot-Savart law, and magnetic forces
\section{Lecture 14 -- Magnetostatics, Biot-Savart Law, and Magnetic Forces}\label{sec:lec14}

\begin{center}
\begin{tabular}{ll}
  \textbf{Date:}\quad 24/05/2026 & \\
\end{tabular}
\end{center}
\hrule\vspace{1.5em}

\subsection*{Overview}

This lecture starts a new static sector of Maxwell theory.  In electrostatics we
isolated the time-reversal invariant sector: no currents, no magnetic field, and
only the scalar potential $\phi$ remained.  Magnetostatics is different.  The fields
are still time independent, but the sources are steady currents.  The main results are:
\begin{itemize}
  \item magnetostatics describes neutral systems with steady currents, not static
        point charges sitting at rest;
  \item the equations reduce to
        $\divg\bB=0$ and $\nabla\times\bB = \frac{4\pi}{c}\Jvec$;
  \item writing $\bB=\nabla\times\Avec$ and choosing the Coulomb gauge
        $\divg\Avec=0$ turns the problem into three Poisson equations;
  \item solving those Poisson equations gives the vector potential
        $\Avec(\br)=\frac{1}{c}\int \Jvec(\br')/|\br-\br'|\,d^3r'$;
  \item taking a curl gives the Biot-Savart law;
  \item for thin wires the Biot-Savart law reduces to the familiar line integral;
  \item the Lorentz force between two current loops can be written as a double line
        integral, and in the static limit Newton's third law is recovered.
\end{itemize}

The guiding idea is the same as in electrostatics: choose variables and gauge so that
Maxwell's equations become Poisson equations.  The difference is that the source is
now a vector field $\Jvec$, so the natural potential is the vector potential $\Avec$.

%------------------------------------------------------------
\subsection{What Magnetostatics Means}\label{subsec:lec14-meaning}

\subsubsection*{Physical Motivation}
The word ``static'' can be misleading.  In magnetostatics charges are moving:
otherwise there would be no current.  What is static is the \emph{macroscopic
configuration}.  A wire carrying a steady current contains moving electrons, but
the positive ionic lattice makes the wire approximately neutral, and the current
distribution does not change with time.

This is why magnetostatics is not obtained by demanding invariance under time
reversal.  Under time reversal,
\begin{equation}\label{eq:lec14-time-reversal}
  t\mapsto -t,\qquad
  \Jvec\mapsto -\Jvec,\qquad
  \bB\mapsto -\bB,\qquad
  \bE\mapsto \bE .
\end{equation}
A current-carrying solution is therefore generally \emph{not} time-reversal
invariant.  It is invariant only under continuous time translations
$t\mapsto t+a$.

\dfn[dfn:lec14-magnetostatics]{Magnetostatic Configuration}{
  A magnetostatic configuration is a time-independent current distribution
  $\Jvec(\br)$ satisfying charge conservation, with no time-dependent charge
  accumulation.  In the simplest neutral sector used in this lecture,
  \begin{equation}\label{eq:lec14-magnetostatic-assumptions}
    \frac{\partial}{\partial t}=0,\qquad \rho=0,\qquad \bE=0,\qquad \Jvec\neq 0 .
  \end{equation}
}

\textit{Significance:} Magnetostatics is a sector of solutions, not a symmetry of
the full theory.  A concrete current pattern usually breaks time reversal even though
Maxwell's equations themselves have the appropriate symmetry structure.

\begin{remark}
  A physical wire is neutral overall: the moving electron charge is balanced by the
  positive lattice.  The current is the charge flow through a cross-section per unit
  time, but the total charge density can still vanish macroscopically.
\end{remark}

%------------------------------------------------------------
\subsection{Maxwell Equations in the Magnetostatic Limit}\label{subsec:lec14-maxwell-limit}

\subsubsection*{Mathematical Setup}
In Gaussian units, Maxwell's equations are
\begin{align}
  \divg\bB &= 0, \label{eq:lec14-maxwell-divB}\\
  \nabla\times\bE + \frac{1}{c}\frac{\partial\bB}{\partial t} &= 0, \label{eq:lec14-maxwell-faraday}\\
  \divg\bE &= 4\pi\rho, \label{eq:lec14-maxwell-gaussE}\\
  \nabla\times\bB &= \frac{4\pi}{c}\Jvec
    + \frac{1}{c}\frac{\partial\bE}{\partial t}. \label{eq:lec14-maxwell-ampere}
\end{align}
With the assumptions \eqref{eq:lec14-magnetostatic-assumptions}, Faraday's law and
Gauss's law for $\bE$ become trivial in the neutral sector, and Ampere-Maxwell loses
the displacement-current term.  We are left with
\begin{equation}\label{eq:lec14-magnetostatic-equations}
  \boxed{
    \divg\bB = 0,\qquad
    \nabla\times\bB = \frac{4\pi}{c}\Jvec .
  }
\end{equation}
\textit{Significance:} Magnetostatics is the magnetic analogue of Poisson theory:
given a steady current $\Jvec$, solve for the magnetic field $\bB$.

Charge conservation also simplifies.  The continuity equation is
\begin{equation}\label{eq:lec14-continuity-full}
  \frac{\partial\rho}{\partial t}+\divg\Jvec=0.
\end{equation}
For time-independent sources, $\partial_t\rho=0$, hence
\begin{equation}\label{eq:lec14-current-conservation}
  \boxed{\divg\Jvec=0.}
\end{equation}
\textit{Significance:} Steady currents cannot begin or end inside the domain; current
lines either close on themselves or enter and leave through boundaries in a balanced way.

%------------------------------------------------------------
\subsection{Introducing the Vector Potential}\label{subsec:lec14-vector-potential}

\subsubsection*{Physical Motivation}
The equation $\divg\bB=0$ says that magnetic field lines have no sources or sinks.
Mathematically, a divergence-free vector field can be represented locally as the curl
of another vector field.  We therefore write
\begin{equation}\label{eq:lec14-B-curl-A}
  \boxed{\bB=\nabla\times\Avec.}
\end{equation}
\textit{Significance:} This automatically solves $\divg\bB=0$, because the divergence
of a curl vanishes identically.

Substituting \eqref{eq:lec14-B-curl-A} into Ampere's law gives
\begin{equation}\label{eq:lec14-curlcurlA-equation}
  \nabla\times(\nabla\times\Avec)=\frac{4\pi}{c}\Jvec .
\end{equation}
To simplify the left-hand side, use the standard vector identity
\begin{equation}\label{eq:lec14-curlcurl-identity}
  \boxed{
    \nabla\times(\nabla\times\Avec)
    = \nabla(\divg\Avec)-\lapl\Avec .
  }
\end{equation}
\textit{Significance:} The only obstruction to a clean Poisson equation is the
gradient term $\nabla(\divg\Avec)$.

\subsubsection*{Index Derivation of the Identity}
Let $\epsilon_{ijk}$ be the three-dimensional Levi-Civita symbol, and use repeated-index
summation.  The $i$th component of $\nabla\times(\nabla\times\Avec)$ is
\begin{align}
  \bigl[\nabla\times(\nabla\times\Avec)\bigr]_i
  &= \epsilon_{ijk}\partial_j(\epsilon_{k\ell m}\partial_\ell A_m) \notag\\
  &= \epsilon_{ijk}\epsilon_{k\ell m}\partial_j\partial_\ell A_m. \label{eq:lec14-curlcurl-index-start}
\end{align}
Using
\begin{equation}\label{eq:lec14-epsilon-contraction}
  \epsilon_{ijk}\epsilon_{k\ell m}
  = \delta_{i\ell}\delta_{jm}-\delta_{im}\delta_{j\ell},
\end{equation}
we get
\begin{align}
  \bigl[\nabla\times(\nabla\times\Avec)\bigr]_i
  &= (\delta_{i\ell}\delta_{jm}-\delta_{im}\delta_{j\ell})
     \partial_j\partial_\ell A_m \notag\\
  &= \partial_i(\partial_j A_j)-\partial_j\partial_j A_i \notag\\
  &= \partial_i(\partial_j A_j)-\lapl A_i. \label{eq:lec14-curlcurl-index-final}
\end{align}
This is exactly \eqref{eq:lec14-curlcurl-identity}.

%------------------------------------------------------------
\subsection{Coulomb Gauge and the Poisson Equation for \texorpdfstring{$\Avec$}{A}}\label{subsec:lec14-coulomb-gauge}

\subsubsection*{Conceptual Background}
The vector potential is not unique.  If $f(\br)$ is any smooth scalar function, then
\begin{equation}\label{eq:lec14-gauge-transformation}
  \Avec\mapsto \Avec'=\Avec+\nabla f
\end{equation}
does not change the magnetic field:
\begin{equation}\label{eq:lec14-gauge-B-invariant}
  \nabla\times\Avec'
  =\nabla\times\Avec+\nabla\times(\nabla f)
  =\nabla\times\Avec.
\end{equation}
This freedom is not a nuisance; it is a tool.  We use it to choose
\begin{equation}\label{eq:lec14-coulomb-gauge}
  \boxed{\divg\Avec=0.}
\end{equation}
\textit{Significance:} The Coulomb gauge removes the gradient term in
\eqref{eq:lec14-curlcurl-identity}, reducing Ampere's law to a vector Poisson equation.

Why is the gauge choice allowed?  Suppose a preliminary potential $\Avec$ does not obey
\eqref{eq:lec14-coulomb-gauge}.  We want $\Avec'=\Avec+\nabla f$ to obey it:
\begin{equation}\label{eq:lec14-gauge-condition-derivation}
  0=\divg\Avec'=\divg\Avec+\lapl f.
\end{equation}
Therefore $f$ must solve
\begin{equation}\label{eq:lec14-gauge-poisson}
  \boxed{\lapl f=-\divg\Avec.}
\end{equation}
\textit{Significance:} Choosing the Coulomb gauge requires solving one scalar Poisson
equation, which is exactly the type of problem already solved in electrostatics.

In Coulomb gauge, \eqref{eq:lec14-curlcurlA-equation} becomes
\begin{equation}\label{eq:lec14-vector-poisson}
  \boxed{-\lapl\Avec=\frac{4\pi}{c}\Jvec.}
\end{equation}
\textit{Significance:} Magnetostatics has been reduced to three copies of Poisson's
equation, one for each Cartesian component of $\Avec$.

For localized currents in free space, the solution is
\begin{equation}\label{eq:lec14-vector-potential-solution}
  \boxed{
    \Avec(\br)=\frac{1}{c}\int
    \frac{\Jvec(\br')}{|\br-\br'|}\,d^3r' .
  }
\end{equation}
\textit{Significance:} The vector potential of a steady current distribution is the
current-density analogue of the electrostatic potential
$\phi(\br)=\int \rho(\br')/|\br-\br'|\,d^3r'$.

\subsubsection*{Consistency Check: \texorpdfstring{$\divg\Avec=0$}{div A = 0}}
Although \eqref{eq:lec14-vector-potential-solution} was obtained in Coulomb gauge, it is
worth checking explicitly that it satisfies the gauge condition.  Taking the divergence
with respect to $\br$,
\begin{align}
  \divg\Avec(\br)
  &= \frac{1}{c}\int
     \Jvec(\br')\cdot\nabla_{\br}\frac{1}{|\br-\br'|}\,d^3r' . \label{eq:lec14-divA-start}
\end{align}
Since
\begin{equation}\label{eq:lec14-grad-switch}
  \nabla_{\br}\frac{1}{|\br-\br'|}
  =-\nabla_{\br'}\frac{1}{|\br-\br'|},
\end{equation}
we can write
\begin{align}
  \divg\Avec(\br)
  &= -\frac{1}{c}\int
     \Jvec(\br')\cdot\nabla_{\br'}\frac{1}{|\br-\br'|}\,d^3r' \notag\\
  &= -\frac{1}{c}\int
     \nabla_{\br'}\cdot
     \left(\frac{\Jvec(\br')}{|\br-\br'|}\right)d^3r'
     +\frac{1}{c}\int
     \frac{\nabla_{\br'}\cdot\Jvec(\br')}{|\br-\br'|}\,d^3r'. \label{eq:lec14-divA-product}
\end{align}
The second integral vanishes by $\divg\Jvec=0$.  The first integral becomes a surface
flux at infinity.  For localized currents, $\Jvec=0$ on a sufficiently large surface,
so the flux vanishes.  Hence
\begin{equation}\label{eq:lec14-divA-zero}
  \boxed{\divg\Avec=0.}
\end{equation}
\textit{Significance:} The current-conservation condition is exactly what makes the
Coulomb-gauge solution self-consistent.

%------------------------------------------------------------
\subsection{The Biot-Savart Law}\label{subsec:lec14-biot-savart}

\subsubsection*{Mathematical Setup}
The magnetic field is the curl of \eqref{eq:lec14-vector-potential-solution}:
\begin{align}
  \bB(\br)
  &= \nabla_{\br}\times\Avec(\br) \notag\\
  &= \frac{1}{c}\int
     \nabla_{\br}\times
     \left(\frac{\Jvec(\br')}{|\br-\br'|}\right)d^3r'. \label{eq:lec14-B-curl-integral}
\end{align}
Inside the curl, $\Jvec(\br')$ is independent of $\br$, so
\begin{equation}\label{eq:lec14-curl-fJ}
  \nabla_{\br}\times\left(f(\br)\Jvec(\br')\right)
  = \nabla_{\br}f(\br)\times\Jvec(\br').
\end{equation}
With
\begin{equation}\label{eq:lec14-grad-green}
  \nabla_{\br}\frac{1}{|\br-\br'|}
  = -\frac{\br-\br'}{|\br-\br'|^3},
\end{equation}
we obtain
\begin{equation}\label{eq:lec14-biot-savart-volume}
  \boxed{
    \bB(\br)=\frac{1}{c}\int
    \frac{\Jvec(\br')\times(\br-\br')}{|\br-\br'|^3}\,d^3r' .
  }
\end{equation}
\textit{Significance:} The magnetic field at $\br$ is built by summing the transverse
contribution of every current element; only the part of the current perpendicular to
the line of sight contributes.

\begin{figure}[h]
\centering
\begin{tikzpicture}[>=Stealth, thick, scale=1.05]
  \draw[blue!60!black, thick]
    plot[smooth cycle, tension=0.8]
    coordinates {(-2,0.2)(-1.2,1.1)(0.3,0.9)(1.2,0.1)(0.7,-0.9)(-0.9,-1.0)};
  \draw[->, blue!60!black] (-1.7,0.55) -- (-1.42,0.82);
  \node[blue!60!black] at (-1.9,1.15) {$C$};
  \fill[red!70!black] (-0.2,0.85) circle (2pt) node[above] {$\br'$};
  \draw[->, red!70!black] (-0.2,0.85) -- (0.45,0.72)
    node[midway, above] {$d\bvec{\ell}'$};
  \fill[black] (2.7,1.4) circle (2pt) node[above] {$\br$};
  \draw[->, gray!80!black] (-0.2,0.85) -- (2.7,1.4)
    node[midway, above] {$\br-\br'$};
  \draw[->, green!50!black] (2.7,1.4) -- (2.7,2.05)
    node[right] {$d\bB$};
  \node[align=center] at (-0.35,-1.55)
    {\footnotesize thin wire: cross-section\\[-1pt]\footnotesize small compared with distances};
\end{tikzpicture}
\caption{Geometry of the thin-wire Biot-Savart law.  A current element
$d\bvec{\ell}'$ at source point $\br'$ contributes a field at observation point $\br$
in the direction $d\bvec{\ell}'\times(\br-\br')$.}
\label{fig:lec14-biot-savart-geometry}
\end{figure}

\subsubsection*{Thin-Wire Limit}
For a wire whose cross-section is small compared with all distances of interest, write
the volume element as $d^3r'=dS'\,d\ell'$ and use
\begin{equation}\label{eq:lec14-current-line-replacement}
  \Jvec(\br')\,dS'\,d\ell' = I\,d\bvec{\ell}' .
\end{equation}
Then \eqref{eq:lec14-biot-savart-volume} becomes
\begin{equation}\label{eq:lec14-biot-savart-line}
  \boxed{
    \bB(\br)=\frac{I}{c}\oint_C
    \frac{d\bvec{\ell}'\times(\br-\br')}{|\br-\br'|^3}.
  }
\end{equation}
\textit{Significance:} The familiar Biot-Savart line integral is a controlled
approximation to the volume-current formula, valid when the wire is thin on the
geometric scales being probed.

\begin{remark}
  A truly open steady wire is impossible in an isolated magnetostatic problem: if
  current entered an endpoint and did not leave, charge would accumulate there.
  That is why the ideal line current is naturally a closed loop, or part of a larger
  closed circuit.
\end{remark}

%------------------------------------------------------------
\subsection{Force Between Two Current Loops}\label{subsec:lec14-loop-force}

\subsubsection*{Physical Motivation}
Once $\bB$ is known, the force on moving charges is the magnetic part of the Lorentz
force.  For two current loops, the field of loop $2$ exerts a force on the charges
moving in loop $1$.  The self-field of loop $1$ gives internal forces and is not part
of the external force between the loops.

Let loop $1$ carry current $I_1$ and loop $2$ carry current $I_2$.  For source points
$\br_1\in C_1$ and $\br_2\in C_2$, define
\begin{equation}\label{eq:lec14-r12-definition}
  \br_{12}=\br_1-\br_2,\qquad r_{12}=|\br_{12}|.
\end{equation}
The force on a small element $d\bvec{\ell}_1$ of loop $1$ due to the magnetic field
created by loop $2$ is
\begin{equation}\label{eq:lec14-element-force}
  d\bvec{F}_{2\to1}
  =\frac{I_1}{c}\,d\bvec{\ell}_1\times\bB_2(\br_1).
\end{equation}
Substituting Biot-Savart for $\bB_2$,
\begin{equation}\label{eq:lec14-loop-force-unsimplified}
  \bvec{F}_{2\to1}
  =\frac{I_1I_2}{c^2}
  \oint_{C_1}\oint_{C_2}
  \frac{d\bvec{\ell}_1\times\left(d\bvec{\ell}_2\times\br_{12}\right)}
       {r_{12}^3}.
\end{equation}
This is correct, but the action-reaction symmetry is not yet transparent.

\subsubsection*{Simplifying the Double Integral}
Use the vector triple-product identity
\begin{equation}\label{eq:lec14-triple-product}
  \bvec{a}\times(\bvec{b}\times\bvec{c})
  = \bvec{b}(\bvec{a}\cdot\bvec{c})
    - \bvec{c}(\bvec{a}\cdot\bvec{b}).
\end{equation}
With $\bvec{a}=d\bvec{\ell}_1$, $\bvec{b}=d\bvec{\ell}_2$, and
$\bvec{c}=\br_{12}$,
\begin{equation}\label{eq:lec14-loop-force-split}
  \frac{d\bvec{\ell}_1\times(d\bvec{\ell}_2\times\br_{12})}{r_{12}^3}
  =
  d\bvec{\ell}_2\,
  \frac{d\bvec{\ell}_1\cdot\br_{12}}{r_{12}^3}
  -
  \br_{12}\,
  \frac{d\bvec{\ell}_1\cdot d\bvec{\ell}_2}{r_{12}^3}.
\end{equation}
The first term vanishes after integration around the closed loop $C_1$.  Indeed,
for fixed $\br_2$,
\begin{equation}\label{eq:lec14-gradient-line-integral}
  \frac{\br_{12}}{r_{12}^3}
  = -\nabla_{\br_1}\frac{1}{r_{12}},
\end{equation}
so
\begin{equation}\label{eq:lec14-closed-gradient-zero}
  \oint_{C_1} d\bvec{\ell}_1\cdot\frac{\br_{12}}{r_{12}^3}
  =
  -\oint_{C_1} d\bvec{\ell}_1\cdot\nabla_{\br_1}\frac{1}{r_{12}}
  =0.
\end{equation}
The last equality says that the integral of a gradient around a closed path is zero.
Therefore
\begin{equation}\label{eq:lec14-loop-force-final}
  \boxed{
  \bvec{F}_{2\to1}
  =-\frac{I_1I_2}{c^2}
  \oint_{C_1}\oint_{C_2}
  \frac{(d\bvec{\ell}_1\cdot d\bvec{\ell}_2)\,\br_{12}}{r_{12}^3}.
  }
\end{equation}
\textit{Significance:} In the static limit the force law becomes explicitly compatible
with Newton's third law, because interchanging $1$ and $2$ sends $\br_{12}\mapsto-\br_{12}$.

Thus
\begin{equation}\label{eq:lec14-newton-third-law}
  \boxed{\bvec{F}_{2\to1}=-\bvec{F}_{1\to2}.}
\end{equation}
\textit{Significance:} Full electrodynamics does not generally obey a particle-only
action-reaction law because fields carry momentum, but the magnetostatic closed-loop
limit restores the expected equal-and-opposite mechanical forces.

%------------------------------------------------------------
\subsection{Parallel Wires as a Local Limit}\label{subsec:lec14-parallel-wires}

\subsubsection*{Physical Motivation}
The usual force between two long parallel wires is an idealization.  Infinite wires do
not exist physically, but two large nearby loops look locally like parallel straight
wires if their radius of curvature is much larger than their separation.

\begin{figure}[h]
\centering
\begin{tikzpicture}[>=Stealth, thick, scale=1.0]
  \draw[->, gray] (-0.5,0) -- (5.2,0) node[right] {$x$};
  \draw[->, gray] (0,-0.5) -- (0,3.0) node[above] {$y$};
  \draw[blue!65!black, line width=1.2pt] (0.4,0.7) -- (4.8,0.7);
  \draw[blue!65!black, line width=1.2pt] (0.4,2.1) -- (4.8,2.1);
  \draw[->, red!70!black] (1.0,0.7) -- (1.7,0.7) node[midway, below] {$I_1$};
  \draw[->, red!70!black] (1.0,2.1) -- (1.7,2.1) node[midway, above] {$I_2$};
  \draw[<->] (4.4,0.7) -- (4.4,2.1) node[midway, right] {$d$};
  \fill (2.1,0.7) circle (2pt) node[below] {$\br_1$};
  \fill (3.5,2.1) circle (2pt) node[above] {$\br_2$};
  \draw[->, gray!80!black] (3.5,2.1) -- (2.1,0.7)
    node[midway, left] {$\br_{12}$};
  \draw[->, green!50!black] (2.1,0.7) -- (2.1,1.25)
    node[right] {$\bvec{F}_{2\to1}$};
  \node at (2.6,0.25) {\footnotesize same-direction currents attract};
\end{tikzpicture}
\caption{Parallel-wire limit of the loop-force formula.  For currents in the same
direction, $d\bvec{\ell}_1\cdot d\bvec{\ell}_2>0$, and the force on the lower wire points
toward the upper wire.}
\label{fig:lec14-parallel-wires}
\end{figure}

For two long parallel wires separated by distance $d$, the force is proportional to the
length $L$ of the wire segment considered.  Therefore the meaningful quantity is the
force per unit length:
\begin{equation}\label{eq:lec14-parallel-wire-force}
  \boxed{
    \frac{F}{L}=\frac{2I_1I_2}{c^2 d}
  }
\end{equation}
with attraction for currents in the same direction and repulsion for currents in
opposite directions.

\textit{Significance:} The sign is not mysterious: the factor
$d\bvec{\ell}_1\cdot d\bvec{\ell}_2$ knows whether the currents are parallel or
antiparallel, and the vector $\br_{12}$ points from the source element to the element
being pushed.

%------------------------------------------------------------
\subsection{Summary}\label{subsec:lec14-summary}

\begin{itemize}
  \item Magnetostatics means time-independent fields with steady currents, not
        time-reversal invariant solutions.
  \item The equations reduce to $\divg\bB=0$ and
        $\nabla\times\bB=\frac{4\pi}{c}\Jvec$, with $\divg\Jvec=0$.
  \item Since $\divg\bB=0$, write $\bB=\nabla\times\Avec$.
  \item Gauge freedom $\Avec\mapsto\Avec+\nabla f$ allows the Coulomb gauge
        $\divg\Avec=0$.
  \item In Coulomb gauge, $\Avec$ satisfies
        $-\lapl\Avec=\frac{4\pi}{c}\Jvec$ and therefore
        $\Avec(\br)=\frac{1}{c}\int \Jvec(\br')/|\br-\br'|\,d^3r'$.
  \item Taking a curl gives the Biot-Savart law
        $\bB(\br)=\frac{1}{c}\int
        \Jvec(\br')\times(\br-\br')/|\br-\br'|^3\,d^3r'$.
  \item For a thin closed wire,
        $\bB(\br)=\frac{I}{c}\oint d\bvec{\ell}'\times(\br-\br')/|\br-\br'|^3$.
  \item The force between two current loops is
        $\bvec{F}_{2\to1}=-(I_1I_2/c^2)\oint\!\oint
        (d\bvec{\ell}_1\cdot d\bvec{\ell}_2)\br_{12}/r_{12}^3$.
  \item For long parallel wires,
        $F/L=2I_1I_2/(c^2d)$, attractive for currents in the same direction.
\end{itemize}
